\completa{30}{Sergio}{5e.1: comparar cada mes con el mismo mes del año anterior}
\senalar{Leer las cuatro cajas; señalar la doble condición del JOIN y la fórmula de variación.}
\decir{La serie mensual permite preguntar cómo cambia cada mes frente al mismo mes del año anterior y cuántos meses suben o bajan. Es la primera consulta de cinco e sobre vistas temporales. Comparar enero con el enero anterior mantiene el mes de referencia; comparar enero con diciembre respondería otra pregunta.}
\decir{El primer cuadro dice «SUM por año y mes; dos copias del agregado». Las dos subconsultas calculan el mismo acumulado mensual y reciben los alias a y b. La flecha conduce a «JOIN consigo misma; mismo mes, año anterior». Esto se llama self join, o unión de una tabla consigo misma. La condición ON exige a punto mes igual a b punto mes y a punto año igual a b punto año más uno. A representa el periodo actual y b el anterior.}
\decir{La siguiente flecha llega a «cien por actual menos previo, dividido entre previo». Primero restamos los conteos, luego dividimos entre el anterior y multiplicamos por cien. CAST AS DOUBLE permite el cálculo decimal y ROUND lo redondea a dos posiciones. La última flecha conduce a «sesenta comparaciones; veintiuna suben, treinta y nueve bajan». Hay cinco años comparables, de 2020 a 2024, por doce meses. El año 2019 sirve de base inicial y no tiene un 2018 en el archivo para compararse.}
\decir{En la salida, enero de 2020 tiene mil ciento sesenta y cinco frente a dos mil ochenta y seis: menos cuarenta y cuatro coma quince por ciento. Febrero tiene dos mil setecientas veinticinco frente a dos mil trescientas veinticuatro: más diecisiete coma veinticinco por ciento. Abril tiene seiscientas cuarenta y tres frente a mil novecientas noventa y seis: menos sesenta y siete coma setenta y nueve por ciento. Son cambios de conteos frente al año anterior, no restas de participaciones. La consulta describe la variación; no identifica su causa. Pasaremos ahora a participaciones territoriales dentro de cada periodo.}

\completa{31}{Sergio}{5e.3: tres tablas, dos JOIN y una ventana por periodo}
\senalar{Seguir las cuatro cajas y marcar el PARTITION BY periodo. La tabla de salida se lee por columnas temporales.}
\decir{Esta consulta de cinco e pregunta: ¿cuál es el peso de cada región dentro de cada periodo tras unir geografía y tiempo? Nace de la posibilidad de relacionar las denuncias con dos catálogos: ubicación por ubigeo y periodo por año. Queremos comparar la composición regional de cada periodo, además de sus acumulados.}
\decir{El primer cuadro dice «Denuncias y dimensiones». La flecha conduce a «Dos JOIN: ubigeo y año». Denuncias se llama d; geografía, g; tiempo, t. Unimos d punto ubigeo con g punto ubigeo y d punto anio con t punto anio. Son tres tablas relacionadas mediante dos JOIN, aunque el comentario original del script diga unión triple. Cada alias identifica de qué tabla procede una columna.}
\decir{La siguiente flecha lleva a «GROUP BY región, periodo; ventana por periodo». Primero sumamos cantidad para cada combinación. Después, SUM de SUM de cantidad OVER PARTITION BY periodo obtiene el denominador de ese periodo. Partition By separa los grupos de la ventana según el periodo; no cambia aquí el grano de la fila de salida ni describe particiones físicas de un archivo. CAST AS DOUBLE permite decimales, ROUND redondea y AVG calcula el promedio por registro que también incluye el código.}
\decir{La última flecha conduce a «nueve combinaciones; total ciento veintiún mil cuatrocientas veintisiete». Son tres regiones por tres periodos. La tabla muestra Costa con sesenta y siete coma ochenta y dos por ciento en 2019, sesenta y tres coma setenta en 2020 a 2021 y sesenta y cinco coma ochenta y ocho en 2022 a 2024. Sierra presenta veintiséis coma quince, veintinueve coma sesenta y veintiocho coma treinta. Selva presenta seis coma cero dos, seis coma sesenta y nueve y cinco coma ochenta y dos.}
\decir{Cada columna temporal se interpreta respecto de su propio total; las participaciones suman aproximadamente cien por ciento por el redondeo. Costa conserva el mayor peso. A diferencia de cinco e dos, donde OVER vacío usa todos los grupos del país, aquí la ventana se restringe al periodo. La suma de las nueve combinaciones conserva el total original. Revisaremos ahora filtros y ordenamientos sobre nacionalidades.}

\completa{32}{Sergio}{5f.2: ordenar acumulados no equivale a medir crecimiento}
\senalar{Leer los cuatro cuadros y la doble regla de ORDER BY.}
\decir{El segundo ejemplo de cinco f pregunta: ¿qué pares de nacionalidad extranjera y año acumulan más personas desde 2022, con al menos diez por grupo? La pregunta se construye con NACIONALIDAD, AÑO y CANTIDAD. Buscamos ordenar volúmenes anuales de nacionalidades declaradas.}
\decir{El primer cuadro dice «WHERE: extranjeros desde 2022». NOT IN excluye PERÚ y NO ESPECIFICA; la condición año mayor o igual a 2022 conserva el periodo final. La flecha conduce a «GROUP BY: nacionalidad y año». Cada pareja recibe una suma de cantidad, un count de registros y un promedio con AVG. ROUND presenta ese promedio con tres decimales.}
\decir{La siguiente flecha llega a «HAVING SUM mayor o igual a diez». Filtramos grupos, no filas originales. La última conduce a «doce pares; personas descendente, año ascendente». El primer criterio ordena de mayor a menor acumulado. Si dos grupos empatan en personas, el segundo coloca primero el año menor. ASC significa ascendente y DESC, descendente. No estamos ordenando principalmente por año.}
\decir{La salida muestra Venezuela en 2023 con cuatrocientas treinta personas; en 2024 con cuatrocientas veintiocho; y en 2022 con cuatrocientas cuatro. Colombia en 2023 aparece con veinticuatro. Venezuela ocupa los tres primeros puestos. Aunque el comentario original del script mencione crecimiento, la consulta ejecutada ordena acumulados y no calcula una tasa de crecimiento: no resta un año anterior ni divide entre su valor. La pregunta de esta ficha corresponde a lo que efectivamente calcula el código. La última consulta filtrada se centrará en estudiantes.}

\completa{33}{Sergio}{5f.3: qué incluye el subconjunto de estudiantes}
\senalar{Seguir las cuatro cajas; señalar los tres rangos de edad y que GROUP BY mes reúne los años.}
\decir{La ocupación y la edad permiten preguntar: ¿en qué meses se acumulan más estudiantes reportados de seis a veintinueve años? Esta consulta completa cinco f. Su universo incluye solo ocupación declarada ESTUDIANTE y los rangos seis a once, doce a diecisiete y dieciocho a veintinueve.}
\decir{El primer cuadro dice «WHERE: estudiante y edad». El operador IN comprueba si rango edad pertenece a esa lista; AND exige además que la ocupación coincida. La flecha lleva a «GROUP BY mes; SUM de cantidad». Agrupamos por número de mes, reuniendo todos los años del archivo, y sumamos personas. Count asterisco cuenta las filas que aportan al grupo.}
\decir{La siguiente flecha conduce a «HAVING SUM mayor que quinientos». Conservamos los meses cuyo acumulado supera el umbral. La última llega a «doce meses; personas descendente, mes ascendente». Primero ordenamos por personas; en un empate, por el número de mes de menor a mayor.}
\decir{La salida muestra diciembre con dos mil seiscientas setenta y tres personas y dos mil diecinueve registros; noviembre con dos mil quinientas treinta y ocho y mil novecientos sesenta; septiembre con dos mil cuatrocientas diecinueve y mil novecientos tres; abril con mil setecientas setenta y cuatro y mil trescientos setenta y ocho. Diciembre encabeza ese subconjunto. No incluye edades de cero a cinco ni edades desconocidas, y tampoco podemos recuperar estudiantes cuya ocupación no está informada. Por ello describimos estudiantes declarados en esos rangos, sin generalizar a todo el alumnado del país. Terminamos las veinticinco consultas; las siguientes láminas permiten revisar operaciones y evidencias.}
